\begin{abstract}

本文面向算力—电力协同（算电协同）调度，基于某算力网络六个区域、50000 条任务与逐时电力数据，沿"预测—碳感知调度—储能协同—联合优化"四步递进建模并求解。

\textbf{对于问题一，} 完成到达 GPU 需求统计、岭回归预测与可行基础调度：最后 24 个到达小时的预测 WAPE 为 $0.623$，较均值基准 $0.657$ 相对提升约 $5\%$，仍属有限，说明在本文模型族与有限历史长度下可预测增益有限，价值主要在真实到达后的调度；贪心调度总等待 $86$ h、迁移率 $0$、平均时延 $5$ ms、容量利用率约 $40\%$，约束零违例。

\textbf{对于问题二，} 在无储能下以"弃电—外送—购电"三段边际成本引导任务重指派：结算成本由 $-4.198\times10^8$ 元单调降至 $-4.442\times10^8$ 元（省 $2.44\times10^7$ 元），碳排由 $8.85$ 升至最高 $777$ 后回落至 $496$ tCO$_2$，即无储能下成本与碳为真实权衡；逐任务时延 SLA 恒满足。

\textbf{对于问题三，} 对固定负荷建立逐区域储能线性规划（LP-first），成本再降 $3.94\times10^7$ 元至 $-4.592\times10^8$ 元并使碳排归零；$\varepsilon$-约束扫描显示任意碳预算下最优成本不变，即储能同时实现成本最优与零碳。D/E/F 三区把低谷弃电移峰外送，分别增利 $0.97/1.37/1.60\times10^7$ 元量级。

\textbf{对于问题四，} 以四格反事实（F00/F10/F01/F11，incumbent 保证可行域嵌套）分解任务柔性—储能交互：得到 $F_{11}=F_{01}=-4.592\times10^8$ 元、交互项 $I=-2.44\times10^7$ 元，任务迁移的收益在储能存在时边际为零，两者在低成本外送套利上高度替代；新能源缩减扫描下任务迁移独立增益衰减至 $0$ 而储能收益恒正，$\rho=0.4$ 时无储能方案不可行而含储能方案仍可行。

综上，储能以 $3.94\times10^7$ 元换得"成本最优与零碳重合"，任务迁移仅在无储能时再省 $2.44\times10^7$ 元，二者呈明显替代。

\keywords{算电协同；碳感知调度；储能优化；任务迁移；多目标优化}
\end{abstract}